% WEIGHT_ADDENDUM.tex, a starter for the extension in which the institution chooses the weight
% mu (author, 15 October 2026: "do it in an addendum - that would serve both as a starter
% document and might probably get us some scope related sentences in the manuscript").
%
% RULES
%   A statement that rests on a row of skeleton_who_pays.tex cites the row. A statement that follows
%   from a row in one step says so and names the step. Everything else is a conjecture or an
%   open question, labelled W-n, with the world in which it is false. The institution's
%   objective is a construct, to be read as rationalising the rules institutions publish, not
%   as a quantity in the world (measurement map, the mu row). Terms are the manuscript's;
%   checks/verify_docs.py (D4, D5) checks the terms and that every cited row exists.

\documentclass[11pt]{article}
\usepackage[T1]{fontenc}
\usepackage[utf8]{inputenc}
\usepackage{mathptmx}
\usepackage{microtype}
\usepackage[a4paper,margin=2.2cm]{geometry}
\usepackage{amsmath,amssymb}
\usepackage{booktabs}
\usepackage{array}
\usepackage[round]{natbib}
\usepackage{parskip}
\usepackage[hidelinks]{hyperref}

\newcolumntype{P}[1]{>{\raggedright\arraybackslash}p{#1}}

\title{The Institution's Weight: A Starter for the Extension}
\author{Companion note to the manuscript}
\date{15 October 2026}

\begin{document}
\maketitle

\section{The question}

In the manuscript the rule's weight on the base component, $\mu$, is given.
Its motivation is the author's account of information asymmetry and
institutional development (14 October 2026): a credential buys a cleaner
status than the disorganised status that comes from unverified providers, and
that is why status becomes important. An institution allocating positions
cannot see the quality it cares about. What it can read without verifying it is
disorganised: claims, reputations and signals from providers it does not check.
A credential is organised status: standardised, issued under rules, and cheap to
check. The weight $1-\mu$ is how far the rule leans on that organised status,
and $\mu$ how far it still reads the rest. As credentialing becomes organised,
an institution has reason to lean on it more, which is the shift from
disorganised to organised status in the story of institutional development.
Once the credential can be bought, its holders are those who can afford it,
which is what the manuscript studies. The extension asks why an institution
sets $\mu$ where it does. Its answer has two sides. The cost side,
what a higher $\mu$ does to who pays and what they take, is in the manuscript.
The benefit side, what a higher $\mu$ does to the quality of the winner, is not.

\section{What the manuscript already gives}

\begin{center}
\small
\begin{tabular}{@{}P{4.2cm}P{7.8cm}P{3.2cm}@{}}
\toprule
\textbf{Statement} & \textbf{Content} & \textbf{Rests on} \\
\midrule
The gain falls with $\mu$ & With the incumbent scored as an entrant,
  $\Delta(m;\mu)$ does not rise in $\mu$ at any $m$. & M28, K8 \\ \addlinespace
The count falls with $\mu$ & With costs fixed, $k^*$ does not rise in $\mu$. &
  M28, K8 \\ \addlinespace
The entrants' share falls with $\mu$ & The probability that the prize goes to an
  entrant does not rise in $\mu$. & M31, K7 \\ \addlinespace
The benchmark at $\mu\to0$ & $\Delta(m)\to V/(m+2)$ for every $Q$. & M9 \\
  \addlinespace
Entry buys nothing at $\mu=1$ & At $\mu=1$ the score is $r$ for every player,
  so $F=G$, $\varphi=0$ and $\Delta(m)=0$; with $c>0$ and $u$ strictly
  increasing, $\kappa>0$ and no challenger enters. & One step from M1;
  \texttt{no\_entry\_at\_one} \\
\bottomrule
\end{tabular}
\end{center}

So, with the fee, the prize and the wealth profile fixed, a rule that weighs the
base component more draws no more entrants and gives them no larger share of the
prize, and at $\mu=1$ it draws none.

\section{What is missing: the benefit side}

Let $r_W$ be the base component of the winner. The institution's benefit from a
higher $\mu$ is a better $r_W$. Two endpoints follow from the score in one step
each, both proved in \texttt{lean/mathlib/WeightEndpoints.lean}; what lies
between them is settled in \texttt{lean/mathlib/WeightSelection.lean} (M32).

\begin{itemize}
\item \textbf{W-1 (proved).} At $\mu=1$ no challenger enters, every score is
  $r$, and the winner is the player with the highest base component, so $r_W$
  is the maximum of the $Q+1$ draws of $r$ (\texttt{winner\_base\_eq\_sup\_at\_one}).
\item \textbf{W-2 (proved).} At $\mu=0$ every score is $s_ie_i$, so the winner is
  picked by a rule on the bought components alone; with $r$ independent of $s$
  and of a common law, $r_W$ has the law of $r$ (\texttt{score\_at\_zero},
  \texttt{winner\_base\_law}). The incumbent, scored as an entrant, always
  holds the bought component, so the winner is defined.
\end{itemize}

Between the endpoints two forces act. A higher $\mu$ weighs $r$ more in every
score, which helps selection on $r$: with the entrants fixed, the winner's base
component does not fall (M32). It also lowers entry (M28), which moves some
challengers from the entrant law $F$ to the outsider law $G$, and that can hurt:
a common bonus washes out when every challenger holds it, and favours its
holders once some stop buying.

\begin{itemize}
\item \textbf{W-3 (settled, M32).} $\mathbb E[r_W]$ does not fall in $\mu$ with
  the entrants fixed, pointwise and for any laws (\texttt{winner\_base\_mono}).
  With entry responding it can fall. With one challenger, $r$ uniform, a bonus
  $s\equiv1$ and $V=1$, the gain is $\frac12$ at $\mu=\frac12$ and $\frac38$ at
  $\mu=\frac23$; a challenger with $\kappa=\frac7{16}$ enters at the first and
  not at the second, and $\mathbb E[r_W]$ is strictly lower at the second
  (\texttt{selection\_falls\_with\_exit}). The conjecture as stated is false.
\end{itemize}

\section{The institution's objective, as a construct}

A candidate objective is
\[
  W(\mu)=\mathbb E[r_W](\mu)+\alpha\,c\,k^*(\mu)-\nu\,\gamma(\mu),
\]
with $\alpha\ge0$ the value to the institution of a fee collected, $\gamma$ a
cost of reading the unverified part of the score, which rises in $\mu$, and
$\nu\ge0$ its price. In the terms of Section~1, $\nu$ is how disorganised the
uncredentialed information is. Each term is a construct. The recorded datum is the rule an institution
publishes: Chiba City's four points for paid care against the rest of its
index, or the points an Italian contracting authority gives a certificate. The
observer orders published rules by $\mu$ without measuring it (measurement map,
the $\mu$ row), and $W$ is admissible only as an ordering of $\mu$ that
rationalises those choices.

The questions the extension would answer:
\begin{itemize}
\item \textbf{W-4.} When is the optimal $\mu$ interior? At $\mu=1$ the fee term
  is zero (no entrant) and verification is dearest; near $\mu=0$ the winner is
  no better than a draw of $r$ (W-2).
\item \textbf{W-5.} Does a spread of wealth move the optimal $\mu$? The count
  responds to a spread by the sign rule of K2 (M13), so the fee term moves with
  the spread, in a direction set by the marginal entrant against the pivot.
\item \textbf{W-6.} Does better unverified information raise the optimal $\mu$?
  A fall in $\nu$ lowers the cost of weighing the base component. With the
  entrants fixed the benefit side does not fall in $\mu$ (M32), but with entry
  responding it need not rise, so the answer turns on how entry moves.
\item \textbf{W-7.} Institutional development. As credentialing becomes
  organised, so that the credential tracks what the institution values more
  closely, does the optimal $\mu$ fall? In the model the credential's content
  is the law of $s$; a credential that is cleaner in this sense would be one
  whose $s$ is less noisy. That is not yet a primitive the model varies.
\item \textbf{W-8.} When does the institution choose closure, $\mu=0$? Only
  when the fee term and the saving on verification outweigh the loss in $r_W$
  from W-1 to W-2.
\end{itemize}

\section{Literature to read first}

The application search of 14 October 2026 found no paper in which a selector
chooses the weight on a purchasable credential given a cost of verification; that
is a statement about the abstracts read (\texttt{LITERATURE.tex}, the
application search of 14 October 2026), not yet a novelty claim. The nearest models, by what they
share with the extension:
\begin{itemize}
\item \citet{FrankelKartik2022}: a decision maker commits to under-using data
  that agents can manipulate. The weight is chosen; the cost of gaming is
  continuous and there is no fee. \textbf{[A]}.
\item \citet{Hancart2026}: firms choose the accuracy and difficulty of a test
  in competition for candidates. The test is chosen; candidates differ only in
  productivity. Read in Sections 1 to 4. \textbf{[A]}.
\item \citet{LiuGargBorgs2022}, Appendix B.2: with a skill the school cannot
  measure, some randomization of the reward can be optimal for the school.
  Read in full (L32).
\item \citet{BenPorathDekelLipman2014}: a principal allocates one object and
  can verify a report at a cost. Verification is chosen per agent; there is no
  credential. \textbf{[A]}.
\item \citet{StahlStrausz2017WP} and \citet{Lizzeri1999}: certification as a
  market, chosen by sellers, buyers or an intermediary. \textbf{[A]}.
\end{itemize}
Each would need a full read, a record and Lean before the extension cites it
for positioning.

\section{What this gives the manuscript}

One paragraph in Appendix~D of the manuscript (``The institution's weight''),
resting only on the table of Section~2: the model gives the cost side of the
institution's choice of $\mu$ and not the benefit side, and the choice itself is
left to the extension.

\bibliographystyle{plainnat}
\bibliography{refs}

\end{document}
